%! AP Statistics — Unit 7, Lesson 7.2 — Group Activity (Answer Key)
\documentclass[10pt]{article}
\usepackage{apstats-article}
\usepackage{apstats-key}   % pulls in -boxes; do NOT also load -boxes
\newcommand{\TallMath}[1]{$\displaystyle #1\rule[-1.4em]{0pt}{3.2em}$}

\begin{document}

\pageheader{Unit 7, Lesson 7.2}{Group Activity --- Answer Key}
\namepartnerperiod

\vspace{0.1in}
\textbf{Directions:} For each problem, use the full four-step procedure (Parameter,
Conditions, Calculate, Conclude). Show all work.

\begin{tcolorbox}[colback=white, colframe=black!40, title=\textbf{Tier R --- Remediate},
    fonttitle=\bfseries, arc=2mm, left=3mm, right=3mm, top=2mm, bottom=2mm]

\textbf{Sleep Study.} A school counselor wants to estimate the mean nightly sleep
duration for students at her school. A random sample of 25 students yields
$\bar x = 6.8$ hours and $s = 1.2$ hours.

Construct a 90\% confidence interval for the true mean nightly sleep duration.

\textbf{Step 1 --- Parameter:} $\mu = $ \ansline{true mean nightly sleep (hours) for all students at the school}

\textbf{Step 2 --- Conditions:} (conditions are pre-verified for you)
\begin{itemize}
    \item Random: random sample of students. \checkmark
    \item 10\%: $25 \le 10\%$ of all students at the school. \checkmark
    \item Large Sample: $n = 25 < 30$; assume counselor verified no strong skewness or
          outliers in the data. \checkmark
\end{itemize}

\textbf{Step 3 --- Calculate:}

$df = $ \makebox[2cm][l]{\ans{24}} \quad $t^* = $ \makebox[2.5cm][l]{\ans{1.711}} (from $t$-table, 90\% CI)

$SE = s/\sqrt{n} = 1.2/\sqrt{25} = $ \makebox[2.5cm][l]{\ans{0.240}}

$ME = t^* \cdot SE = $ \makebox[4cm][l]{\ans{$1.711(0.240) = 0.411$}}

Interval: $\bar x \pm ME = 6.8 \pm $ \makebox[2cm][l]{\ans{0.411}} $= ($ \makebox[2cm][l]{\ans{6.389}} $,$ \makebox[2cm][l]{\ans{7.211}} $)$ hours

\textbf{Step 4 --- Conclude:} \ansline{We are 90\% confident that the interval (6.39, 7.21) hours captures}\\\ansline{the true mean nightly sleep duration for all students at this school.}

\end{tcolorbox}

\begin{tcolorbox}[colback=white, colframe=black!40, title=\textbf{Tier A --- Approaching Proficiency},
    fonttitle=\bfseries, arc=2mm, left=3mm, right=3mm, top=2mm, bottom=2mm]

\textbf{Commute Times.} A city planner wants to estimate the mean commute time for
workers in a downtown district. A random sample of 36 workers is taken; summary
statistics: $\bar x = 28.4$ min, $s = 9.6$ min.

Construct a 95\% confidence interval for the true mean commute time. Use all four
steps and verify conditions explicitly.

\vspace{0.1in}
\textbf{Step 1:} \ansline{$\mu$ = true mean commute time (min) for all downtown workers.}

\textbf{Step 2 --- Conditions:} \ansline{Random: random sample of 36 workers. \checkmark}\\\ansline{10\%: $36 \le 10\%$ of all downtown workers. \checkmark}\\\ansline{Large Sample: $n = 36 \ge 30$, so the CLT applies. \checkmark}\\\ansline{}

\textbf{Step 3:}

$df = $ \makebox[2cm][l]{\ans{35}} \quad $t^* = $ \makebox[2.5cm][l]{\ans{2.030}} \quad $SE = $ \makebox[2.5cm][l]{\ans{1.60}}

Interval: \makebox[10cm][l]{\ans{$28.4 \pm 2.030(1.60) = 28.4 \pm 3.25$}} $= ($ \makebox[2cm][l]{\ans{25.15}} $,$ \makebox[2cm][l]{\ans{31.65}} $)$ min

\textbf{Step 4:} \ansline{We are 95\% confident that the interval (25.15, 31.65) minutes captures}\\\ansline{the true mean commute time for all workers in the downtown district.}

\vspace{0.05in}
\textbf{Bonus:} How wide must the interval be? How could the planner halve the margin
of error? \ansline{}\\\ansline{Width $= 2(3.25) = 6.50$ min. To halve the ME, quadruple $n$ (36 to 144).}

\end{tcolorbox}

\begin{tcolorbox}[colback=white, colframe=black!40, title=\textbf{Tier E --- Extension},
    fonttitle=\bfseries, arc=2mm, left=3mm, right=3mm, top=2mm, bottom=2mm]

\textbf{Matched-Pairs: Step-Count Intervention.} A fitness researcher measured daily
step counts for 12 participants before and after a 4-week intervention. The differences
(after $-$ before, in thousands of steps) are:

\medskip
\begin{center}
\begin{tabular}{cccccccccccc}
1.2 & 0.8 & $-0.3$ & 2.1 & 1.5 & 0.6 & 1.8 & $-0.1$ & 2.4 & 0.9 & 1.3 & 1.7
\end{tabular}
\end{center}
\medskip

Compute $\bar d$ and $s_d$, verify conditions, and construct a 95\% CI for the mean
change in daily step count. Then interpret the interval and assess whether the data
provide convincing evidence that the intervention increased step counts on average.

$\bar d = $ \makebox[3cm][l]{\ans{1.158}} \quad $s_d = $ \makebox[3cm][l]{\ans{0.784}} \quad $n = $ \makebox[2cm][l]{\ans{12}}

\vspace{0.1in}
\textbf{Step 2 --- Conditions:} \ansline{Random: random sample of participants. \checkmark}\\\ansline{10\%: $12 \le 10\%$ of all such people. \checkmark}\\\ansline{$n = 12 < 30$: differences show no strong skew or outliers. \checkmark}

\textbf{Step 3:} $df = $ \makebox[2cm][l]{\ans{11}} \quad $t^* = $ \makebox[2.5cm][l]{\ans{2.201}}

Interval: $\bar d \pm t^* \cdot s_d/\sqrt{n} = $ \makebox[8cm][l]{\ans{$1.158 \pm 2.201(0.784/\sqrt{12}) = 1.158 \pm 0.498$}} $= ($ \makebox[2cm][l]{\ans{0.660}} $,$ \makebox[2cm][l]{\ans{1.656}} $)$ thousands of steps

\textbf{Step 4 and claim assessment:} \ansline{95\% confident the true mean change is}\\\ansline{between 0.660 and 1.656 thousand steps per day. The whole interval is above 0,}\\\ansline{so there is convincing evidence the intervention increased step counts.}

\end{tcolorbox}

\end{document}
